%%%%%%%%%%%%%%%%%%%%%%%%%%%%%%%%%%%%%%%%%%%%%%%%%%%%%%%%%%%%
%% CHAPTER 1 -- The Idea That a Word Predicts the Next Word
%%%%%%%%%%%%%%%%%%%%%%%%%%%%%%%%%%%%%%%%%%%%%%%%%%%%%%%%%%%%

\chapter{The Idea That a Word Predicts the Next Word}
\label{ntg:ch1}

\scenelabel
\begin{scenestrip}
\sceneitem{The problem}
\sceneitem{How we got here}
\sceneitem{Where we stand}
\sceneitem{What goes wrong}
\end{scenestrip}

\paragraph{The problem.} Consider for a moment what it means to predict the next word. You
are reading this sentence. Your eye has not yet reached its end,
and yet some part of you is already guessing how it will finish.
The guess is not certain. It is not arbitrary either. It is
shaped by everything you have ever read, and by the few words
that came just before. If we wanted to teach a machine to do the
same, we would need to give it a way to assign, to each possible
continuation, a number that captures \emph{how likely} that
continuation is given everything it has seen so far.

That is the entire problem of language modelling, written in a
single sentence. Everything in this guide is, in one way or
another, the elaboration. The Transformer of Chapter~5 is a
particular machinery for assigning those numbers. The training
process of Chapter~6 is what taught the machinery which numbers
to assign. The alignment work of Chapter~9 is what later
intervened to keep the machinery from assigning high numbers to
sentences nobody would want it to produce. The failure modes of
Chapter~13 are, almost without exception, traceable to the gap
between what makes a continuation \emph{likely} and what makes
it \emph{true}.

It is worth being precise about what ``predict the next word''
does and does not mean. It does not mean the machine knows what
the writer intended. It does not mean the machine knows what the
sentence is about. It means, narrowly, that the machine has been
shown enough text to have built up a sense of which sequences of
words have, in the past, been followed by which other words ---
and that it can, when given a new prefix, hand back a list of
candidate continuations with a number attached to each. The
number is meant to capture how often a continuation like this
has, in the machine's experience, followed a prefix like that.
Nothing more. Nothing less.

\paragraph{How we got here.} The idea that one symbol in a sequence depends on the symbols
before it is older than the computer.

In 1913 the Russian mathematician Andrey Markov sat down with a
copy of Pushkin's verse novel \emph{Eugene Onegin} and counted,
by hand, the alternation of vowels and consonants across twenty
thousand letters. He did this for a reason that was, by the
standards of the time, philosophical. A colleague had argued
that the law of large numbers --- the principle that says a
sufficiently long sequence of dice rolls will average out to the
expected value --- required the dice to be independent of one
another. Markov wanted to show that one could relax independence
into a chain of conditional dependence and still do useful
mathematics. He picked Pushkin because the alternation of vowels
and consonants in a poem is obviously not independent (a
consonant is more often followed by a vowel than by another
consonant) and because Pushkin was conveniently long. He
published his result and gave us, almost as a side effect, a
model of language as a sequence of dependencies.

He gave it to us a full century before we had the data or the
hardware to make it useful.

Claude Shannon, in 1948, wrote the paper that founded
information theory. He was working at Bell Labs on the problem of
sending telegraphs and telephone calls efficiently across
imperfect wires. As one of his running examples he chose English
text. Shannon estimated the entropy of English at roughly one bit
per letter --- meaning, in plain language, that each letter of an
English string carries about as much information as a single
flip of a fair coin, after you take into account everything you
already knew from the letters before it. He observed that a model
that captures more of this structure will, by definition, assign
higher probability to natural text than to random text. He gave
us a yardstick by which to measure how good any future language
model would be.

For half a century after Shannon the field crawled. The n-gram
era of speech recognition counted word triples on hard drives
that filled rooms. ``N-gram'' is a piece of jargon that means,
simply, ``a sequence of $n$ consecutive words'' --- so a
trigram is three words, a four-gram is four, and so on. The
recipe was to scan an enormous corpus of text, count how often
each trigram appeared, and use the counts to estimate the
probability that the third word followed the first two. This
worked well enough to enable the first generation of practical
speech recognition, and it worked badly enough that the field
acquired a reputation for being a graveyard of careers. Frederick
Jelinek, working at IBM in the 1980s, gave us the catchphrase
``every time I fire a linguist my speech recognizer's performance
improves.'' It was a remark about the willingness of mere counts
to do the work that elaborate linguistic theory had failed to do.

The neural turn came in 2003. Yoshua Bengio and his coauthors
built a small neural network that learned a continuous
representation of words --- a list of numbers per word, of the
sort we will spend Chapter~4 understanding --- and used those
representations to predict the next one. The network was not
much better than the n-gram models of the day, and the paper
went largely unnoticed outside its specialist community. But
the idea had been planted: instead of treating each word as a
separate symbol whose only relationship to other words was the
co-occurrence count in the corpus, you could give each word a
position in a high-dimensional space and let the geometry of
that space carry meaning.

Tom\'a\v{s} Mikolov and his colleagues, in 2013, simplified
Bengio's idea into something called \emph{word2vec}. The
geometry of the resulting word vectors --- in which ``king minus
man plus woman'' landed near ``queen'' --- caught public
attention in a way that the underlying technology, by then ten
years old, had not. We will spend Chapter~4 unpacking what this
means. For now it is enough to say that the trick was real, the
trick had limits, and the trick has now been absorbed into the
foundations of every model we will discuss.

After 2013 the decade compressed. Sequence-to-sequence
translation arrived in 2014 (Sutskever, Vinyals, Le).
\emph{Attention} arrived in 2015 (Bahdanau, Cho, Bengio). The
Transformer arrived in 2017 (Vaswani and seven coauthors, in a
Google paper titled \emph{Attention Is All You Need}). And then
the GPT line picked up. By the end of 2022 the kind of model
that had been a research curiosity for fifty years was a
consumer product. Almost everyone reading this book has by now
spoken to one.

\paragraph{Where we stand.} The lineage converged on a remarkably simple recipe. Take a
neural network --- we will see in Chapter~3 what one of those is.
Show it a long sequence of words. Ask it to predict, at each
position in the sequence, what the next word is going to be ---
not as a single guess, but as a relative weighting over every
word in its vocabulary. Compare its weighting to the word that
actually came next. Penalise the network in proportion to how
\emph{surprised} it was by the truth: a small penalty if the true
word was high on its list, a large penalty if the true word was
low. Average that penalty over a corpus of billions of words.
Adjust the network, in tiny increments, to reduce the penalty.
Repeat for months, on machines that cost more than airliners.

That is, in cartoon form, what the training of a Large Language
Model is. The technical name for the penalty we just described is
\emph{cross-entropy}. We will not need to compute it, and we will
not write down its formula. What we need to remember is what it
is for. The cross-entropy penalty is the number a training run
spends months trying to make smaller. Every other technical
choice in this book --- the architecture, the data mix, the
optimiser, the learning rate, the size of the model --- is, in
the end, a choice in service of getting that single number down.

A model that has been trained to make this number small has
learned to assign high probability to the kinds of continuation
that actually appeared in its training data, and low probability
to the kinds that did not. That is the entire content of the
training. Everything else --- the apparent reasoning, the
apparent helpfulness, the apparent personality --- is a
consequence of the data the model was trained on, the way the
data was filtered and ordered, and the additional steps we will
describe in Chapter~9 that turned a fluent autocomplete into
something that resembles an assistant.

\paragraph{What goes wrong.} A model trained the way we just described is, by construction,
the most plausible thing the model can say given the prompt. The
most plausible thing is not the same as the true thing. It is
not the same as the most useful thing. It is not the same as
the thing the user wanted, or the thing the user would have
wanted if they had thought about it longer. It is, narrowly, the
continuation that the training corpus would on average have
produced. We give a polite name to the gap between this and the
truth: we call it \emph{hallucination}. Hallucination is not a
bug. It is the training procedure doing exactly what we asked of
it. We will see this point repeated, in different language, in
nearly every chapter that follows.

There is a second consequence worth stating now, because we will
return to it under different names many times. The model does
not have, in any meaningful sense, a sense of how confident it
should be. What it has is a probability distribution over the
next word. Those two things are not the same. A weather forecast
that says ``70\% chance of rain tomorrow'' is meaningful only if
the forecaster is right about the rain roughly seventy times out
of a hundred when they make that call; a forecaster who says
``70\%'' on every cloudy day, regardless of whether it actually
rains, is producing numbers but not producing information. The
technical name for the property of a forecaster being right
about how often they are right is \emph{calibration}, and Large
Language Models are, in general, only weakly calibrated. We can
measure how poorly. We can sometimes correct it. We cannot
eliminate it.

The reader who began this chapter with the intuition that an LLM
``just predicts the next word'' will, by the end of the next few
chapters, have a sturdier intuition: an LLM produces, at each
step, a relative weighting over possible continuations, and the
weighting reflects what its training data said was likely. That
is all it does. The remarkable thing is how much of what feels
like intelligence falls out of that single mechanism. The
worrying thing is how much does not.

\section{The idea, in plain language}

\subsection*{Probability, in the sense the model uses it}

A Large Language Model does not produce a single answer. It
produces, at each step, a list of possible next words and a
number for each one. The number is a relative weight that the
model would put on that word, given everything before it. The
weights add up to one across all the possible words in the
model's vocabulary. So a typical pattern, after the prefix ``the
cat sat on the'', might look something like:

\begin{itemize}
  \item ``mat'' --- weight 0.31
  \item ``floor'' --- weight 0.18
  \item ``couch'' --- weight 0.11
  \item ``windowsill'' --- weight 0.06
  \item \ldots and a long tail of much smaller weights
\end{itemize}

The model, in producing a response, almost never just picks the
top of the list. It usually samples from the list according to
the weights, with a temperature setting --- which we will meet
in Chapter~6 --- that controls how aggressively it favours the
top. Sampling from the list is what gives the model its apparent
creativity, and it is also what gives it its variability: ask
the same prompt twice and you may get two different answers,
because the model rolled a different weighted die.

This is the first piece of intuition worth holding onto. The
model is not making a decision. It is sampling from a
distribution. The distribution is what the model knows. The
single answer you see is one draw from it. The fact that you
see only one draw is the fact that lets the model feel
authoritative when it should feel uncertain.

\subsection*{What the training is, in plain language}

We have already described the training in cartoon form. It is
worth restating, slightly less briefly, because almost
everything in the rest of the guide is downstream of this
description.

The training corpus is a very large body of text. For modern
models it consists of most of the public English-language web
(filtered, deduplicated, and otherwise cleaned), large
collections of books, large collections of source code, and
varying amounts of text in other languages. Estimates of the
size of the training corpora used for the largest current
models run from one to ten trillion words, depending on which
model and how you count. By comparison, an English-speaking
adult who reads voraciously for sixty years might read on the
order of one to two hundred million words --- four to five
orders of magnitude less.

The training procedure walks through this corpus, again and
again, position by position. At each position, the model is
shown the words up to that point and asked to predict the next
one. Its prediction is compared to the actual next word, and a
small adjustment is made to its internal numbers (its
\emph{parameters} --- a modern model has somewhere between
a few billion and a few trillion of them) to make it slightly
less surprised the next time it sees a similar prefix. The
adjustments are tiny. They have to be: a single corpus pass
applies the adjustment trillions of times, and if any one of
them were large the model would oscillate wildly.

The penalty being minimised --- the surprise --- is what
we earlier called cross-entropy. It has several equivalent
descriptions, each of which is useful for a different purpose.
It is the number of bits, on average, the model is using to
encode each word it sees. It is the negative logarithm of the
probability the model assigned to the truth. It is a measure of
how far the model's distribution is from the distribution the
training data appears to have been drawn from. We will not
need any of these formal descriptions, but it is worth knowing
they exist, because every paper you will ever read about
language modelling reports its results in one of them.

\subsection*{Why this works at all}

It is reasonable to ask, on first encounter, why a procedure as
crude as ``predict the next word a few trillion times'' produces
a system that can answer complex questions, write working code,
or hold a passable conversation. The honest answer is that the
field did not, for most of its history, expect that it would.
The early generations of next-word predictors were good at
imitating the surface texture of language and bad at almost
anything else. Even Bengio's 2003 paper, in retrospect a
foundational work, was not received as such at the time.

What changed was scale. As the corpora got larger and the models
got bigger, a series of capabilities began to appear that had
not been explicitly designed in. The model started to be able to
complete arithmetic problems. It started to be able to write
syntactically correct code in languages it had been shown. It
started to be able to translate between languages without
having been trained on translation pairs. None of this was
anticipated. All of it appears to be a consequence of the fact
that a sufficiently large model trained on a sufficiently large
amount of text has, in order to predict next words well, had to
build internal representations that are useful for a great many
tasks the training never asked about. We will return to this in
Chapter~6 under the name \emph{in-context learning}, and in
Chapter~7 under the name \emph{scaling laws}.

The honest answer to the question ``why does this work'' is
that we have a partial empirical understanding and a poor
theoretical one. We know that scale helps. We know roughly
\emph{how much} it helps, in a way we will quantify in Chapter~7.
We do not know, in any deep sense, why next-word prediction
should produce the breadth of behaviour it does. The most
useful posture for a non-technical reader is to take this as
an empirical fact about the universe and to remain
appropriately humble about predicting which capabilities will
or will not survive in the next generation of models.

\section{The engineering consequence}

If everything an LLM does is downstream of next-word prediction,
several practical consequences follow that anyone reasoning
about LLM products should hold in mind.

\paragraph{The model's outputs are samples, not statements.}
Because the model produces a distribution and then samples from
it, the same prompt can produce different answers on different
runs. This is a feature for creative tasks, where it is what
gives the model its apparent imagination. It is a problem for
tasks where reproducibility matters --- a clinical decision
support tool, say, or a legal review --- and the standard fix
is to lower the temperature so that the model almost always
picks the top of its list. Lowering the temperature does not
make the model more correct. It makes it more consistent, which
is not the same thing.

\paragraph{The model's training data is its world.}
The model knows what was in its corpus, in roughly the
proportions in which it appeared. Topics that are rare in the
corpus --- minority languages, specialist subfields, the
contents of paywalled databases, anything published after the
training cut-off --- are correspondingly underrepresented in
what the model knows. A model trained mostly on English-language
text will be measurably worse at Estonian; a model trained
without access to a particular pharmacology textbook will not
have memorised it. This is not a moral failing of the model. It
is a property of the training data, and it can only be
addressed by changing the data, by retrieval (Chapter~10), or
by accepting the limit.

\paragraph{Bigger is not always better, but bigger usually
helps.} The relationship between model size, training data
size, training compute, and resulting capability is not
mysterious. It has been quantified, for the most part, by the
papers we will discuss in Chapter~7, beginning with the work
of Jared Kaplan and his colleagues at OpenAI in 2020 and
sharpened by the work of Jordan Hoffmann and his colleagues at
DeepMind in 2022 (the so-called \emph{Chinchilla} paper). The
relationship is roughly: capability improves as a slow,
predictable function of compute, until the compute is no longer
the binding constraint, after which it stops improving and
sometimes regresses. Many vendor announcements in this field
are best understood as movements along, or away from, that
curve.

\paragraph{The penalty being minimised is the wrong one for many
real tasks.} Cross-entropy on next-word prediction is a proxy
for ``does the model speak like a human''. It is not a proxy
for ``does the model give correct answers'', or ``does the
model behave safely'', or ``does the model help the user
accomplish their goal''. Each of those is a separate property
that has to be added on top, and adding each of them costs
something in the original capability. Chapter~9 is devoted to
the trade-offs of doing this.

\section{What goes wrong, and why}

The failure modes that come from the mechanism described in this
chapter are the ones that will appear, in different forms, in
every later chapter. It is worth stating them once now, in
plain language, so that we can refer back to this section when
they reappear.

\paragraph{Hallucination.} The model produces a fluent,
confident, plausible answer that is in fact wrong. Sometimes
this is a fact (``Python was created in 1989'' --- it was 1991);
sometimes a citation (``as Smith and Jones (2018) showed'' ---
no such paper exists); sometimes a chain of reasoning that goes
wrong in the middle without the model noticing. The mechanism
is the one we have just described. The model assigns weight to
plausible continuations. In the parts of its training corpus
where it has seen a fact stated many times, the most plausible
continuation is the true one. In the parts where it has seen a
fact stated once or never, the most plausible continuation is
whatever the surrounding pattern suggests, which is often
right and is sometimes spectacularly wrong. The model has no
way of telling these two cases apart from the inside. It treats
them identically. This is the central, irreducible failure mode
of the technology, and Chapter~10 (retrieval) and Chapter~13
(the error model) are largely about how the engineering works
around it.

\paragraph{Miscalibration.} The model's expressed confidence
does not reliably match its accuracy. A vanilla pretrained
model is, in fact, reasonably well-calibrated on next-word
prediction itself. The trouble is that the alignment training
of Chapter~9 tends to make the model both more helpful and less
calibrated --- it learns to sound confident because users
prefer confident-sounding answers, and the result is a model
that sometimes asserts things it should hedge and sometimes
hedges things it should assert. This is one of the costs of
the alignment work, and it has not, as of this writing, been
fully solved.

\paragraph{Probability, confidence, and correctness are three
different things.} The model has a probability over the next
word. We tend to read fluent, definitive prose as
\emph{confident}, and we tend to read confident prose as
\emph{correct}. The model's prose is fluent because that is
what its training optimised for. Its confidence in the
colloquial sense is a property of the prose style, not of the
underlying probabilities. Its correctness is a property of the
match between what it said and the world, which is a property
the model has no direct access to. Conflating the three is the
source of more downstream errors than any architectural
choice. We will return to this distinction in Chapter~13 as
the difference between the \emph{meaning} axis and the
\emph{groundedness} axis of the error model.

\paragraph{Sensitivity to small changes in the prompt.} Because
the model's response is one sample from a distribution that
depends on the exact prefix, two prompts that differ in
seemingly trivial ways --- a comma added, a word reordered, a
synonym substituted --- can produce qualitatively different
outputs. The technical name is \emph{prompt brittleness}, and
we will see in Chapter~5 why the architecture makes this not
just possible but inevitable. For the purposes of this chapter,
the important point is that the variability is intrinsic and
that it cannot be removed simply by ``writing better prompts'',
though writing better prompts can certainly reduce it.

\section*{Bridges}
\addcontentsline{toc}{section}{Bridges}

This chapter sets up the foundation for everything that follows.
The notion that a model produces a probability distribution over
next words reappears in Chapter~5 as the output of the
Transformer, in Chapter~6 as the basis of pretraining, in
Chapter~9 as the thing alignment modifies, and in Chapter~13
as the source of the \emph{meaning} and \emph{groundedness} axes
of the error model. The notion that the model is sampling
rather than answering reappears in Chapter~6 as the temperature
control. The notion that the model is calibrated only weakly
reappears in Chapter~9 as a measurable cost of the alignment
work and in Chapter~10 as the reason retrieval, when done well,
helps so much.

This chapter depended on no earlier chapter, and the reader
will not lose anything by stopping here for a day before
continuing. The next chapter --- \emph{How Machines Learn from
Their Mistakes} --- explains, in similarly plain language, the
``adjust the network in tiny increments to reduce the
penalty'' step we have repeatedly waved at.
